\section{At Gethsemane}

The movement by which our obedience is perfected in love is not an escape from concrete action. Gethsemane does not show Christ progressing from loveless duty into charity; it reveals their perfect unity in the most concrete consent of his human life. In the garden he does not deny anguish or pretend that death is naturally desirable. He prays, ``Father, if you are willing, take this cup away from me; still, not my will but yours be done'' (Luke 22:42).

The sentence is sometimes heard as though freedom were the destruction of a human will. Catholic Christology says the opposite. The eternal Son possesses a complete human nature and therefore a real human will. That will is not sinful, fictitious, or absorbed into the divine will; it freely adheres to the Father's saving will in the extremity of suffering (CCC 475). The natural human shrinking from death gives the consent its terrible concreteness. It does not make the consent a contest between a disobedient humanity and an overpowering God.\endnote{CCC 475 receives the dyothelite teaching of the Third Council of Constantinople: Christ's human will follows and is freely subject to, rather than resisting or opposing, his divine will. Benedict XVI's general audience of 1 February 2012 reads Gethsemane as the human will's complete ``yes'' to the divine will and the place where true freedom is realized. The article uses that doctrinal synthesis; it does not speculate about Christ's interior psychology beyond the Gospel and received teaching.}

Gethsemane therefore reveals the summit of created human freedom, not because every free act must take the form of physical abnegation, but because freedom reaches perfection when the real created will adheres without reserve to the true good. Createdness is predicated here of the human will assumed by the eternal Son, never of the divine Person. For most Christians the immediate form of such adherence is not dramatic suffering. It is fidelity to a spouse, care for a child or parent, truthful work, necessary rest, resisted temptation, repaired injustice, prayer kept in dryness, and mercy offered when no emotion carries it.

Christ's consent is more than an example. The Son's human obedience belongs to the unique redemptive offering by which the creature's answer is healed and taken into his. No penance, suffering, or act of worship stands beside it as an equal sacrifice. Christian action is possible within his prior ``yes.''

Justice brought the inquiry to a will that does what is due. In the garden, obedience is revealed from within as filial self-gift. The act remains costly; the command has not disappeared. Yet the final intelligibility of the act is not a court whose claim has at last been satisfied. It is communion with the Father and love for those whom the Son came to save.

The question after vertigo was, \emph{What must I do?} It cannot be skipped. It also cannot be the last question.

\begin{center}
\emph{What must I learn to want?}
\end{center}
